% --- INICIO DE PLANTILLA PERSONALIZADA ---
\documentclass[oneside,11pt]{Latex/Classes/PhDthesisPSnPDF}

% --- Paquetes Originales ---
\usepackage{blindtext}
\usepackage{actuarialangle} 
\usepackage{tabularx}
\usepackage{float}
\usepackage{multirow, array}
\usepackage[usenames,dvipsnames,svgnames,table]{xcolor}
\usepackage{longtable}
\usepackage{latexsym,amsmath,amssymb,amsfonts}
\usepackage{enumitem}
\usepackage{xparse}
\usepackage{booktabs}
\usepackage[round, sort, numbers]{natbib}  
\usepackage{adjustbox}
\usepackage[utf8]{inputenc}
\usepackage{caption}
\usepackage{graphicx}
\usepackage{hyperref}
\usepackage{cleveref}

% Definición de colores
\definecolor{miblue}{RGB}{68, 114, 196}

% --- CAMBIO CRÍTICO 1: Usar \input para macros en el preámbulo ---
% \include da problemas antes del begin document. \input es seguro.
\input{Latex/Macros/MacroFile1}

% --- Definiciones de seguridad (por si la clase falla al cargar alguna variable) ---
\providecommand{\facultad}[1]{\def\lafacultad{#1}}
\providecommand{\escudofacultad}[1]{\def\elescudo{#1}}
\providecommand{\degree}[1]{\def\elgrado{#1}}
\providecommand{\director}[1]{\def\eldirector{#1}}
\providecommand{\lugar}[1]{\def\ellugar{#1}}
\providecommand{\degreedate}[1]{\def\lafecha{#1}}

% --- Configuración para bloques de código de R (Pandoc) ---
\usepackage{color}
\usepackage{fancyvrb}
\newcommand{\VerbBar}{|}
\newcommand{\VERB}{\Verb[commandchars=\\\{\}]}
\DefineVerbatimEnvironment{Highlighting}{Verbatim}{commandchars=\\\{\}}
% Add ',fontsize=\small' for more characters per line
\usepackage{framed}
\definecolor{shadecolor}{RGB}{248,248,248}
\newenvironment{Shaded}{\begin{snugshade}}{\end{snugshade}}
\newcommand{\AlertTok}[1]{\textcolor[rgb]{0.94,0.16,0.16}{#1}}
\newcommand{\AnnotationTok}[1]{\textcolor[rgb]{0.56,0.35,0.01}{\textbf{\textit{#1}}}}
\newcommand{\AttributeTok}[1]{\textcolor[rgb]{0.13,0.29,0.53}{#1}}
\newcommand{\BaseNTok}[1]{\textcolor[rgb]{0.00,0.00,0.81}{#1}}
\newcommand{\BuiltInTok}[1]{#1}
\newcommand{\CharTok}[1]{\textcolor[rgb]{0.31,0.60,0.02}{#1}}
\newcommand{\CommentTok}[1]{\textcolor[rgb]{0.56,0.35,0.01}{\textit{#1}}}
\newcommand{\CommentVarTok}[1]{\textcolor[rgb]{0.56,0.35,0.01}{\textbf{\textit{#1}}}}
\newcommand{\ConstantTok}[1]{\textcolor[rgb]{0.56,0.35,0.01}{#1}}
\newcommand{\ControlFlowTok}[1]{\textcolor[rgb]{0.13,0.29,0.53}{\textbf{#1}}}
\newcommand{\DataTypeTok}[1]{\textcolor[rgb]{0.13,0.29,0.53}{#1}}
\newcommand{\DecValTok}[1]{\textcolor[rgb]{0.00,0.00,0.81}{#1}}
\newcommand{\DocumentationTok}[1]{\textcolor[rgb]{0.56,0.35,0.01}{\textbf{\textit{#1}}}}
\newcommand{\ErrorTok}[1]{\textcolor[rgb]{0.64,0.00,0.00}{\textbf{#1}}}
\newcommand{\ExtensionTok}[1]{#1}
\newcommand{\FloatTok}[1]{\textcolor[rgb]{0.00,0.00,0.81}{#1}}
\newcommand{\FunctionTok}[1]{\textcolor[rgb]{0.13,0.29,0.53}{\textbf{#1}}}
\newcommand{\ImportTok}[1]{#1}
\newcommand{\InformationTok}[1]{\textcolor[rgb]{0.56,0.35,0.01}{\textbf{\textit{#1}}}}
\newcommand{\KeywordTok}[1]{\textcolor[rgb]{0.13,0.29,0.53}{\textbf{#1}}}
\newcommand{\NormalTok}[1]{#1}
\newcommand{\OperatorTok}[1]{\textcolor[rgb]{0.81,0.36,0.00}{\textbf{#1}}}
\newcommand{\OtherTok}[1]{\textcolor[rgb]{0.56,0.35,0.01}{#1}}
\newcommand{\PreprocessorTok}[1]{\textcolor[rgb]{0.56,0.35,0.01}{\textit{#1}}}
\newcommand{\RegionMarkerTok}[1]{#1}
\newcommand{\SpecialCharTok}[1]{\textcolor[rgb]{0.81,0.36,0.00}{\textbf{#1}}}
\newcommand{\SpecialStringTok}[1]{\textcolor[rgb]{0.31,0.60,0.02}{#1}}
\newcommand{\StringTok}[1]{\textcolor[rgb]{0.31,0.60,0.02}{#1}}
\newcommand{\VariableTok}[1]{\textcolor[rgb]{0.00,0.00,0.00}{#1}}
\newcommand{\VerbatimStringTok}[1]{\textcolor[rgb]{0.31,0.60,0.02}{#1}}
\newcommand{\WarningTok}[1]{\textcolor[rgb]{0.56,0.35,0.01}{\textbf{\textit{#1}}}}

%%%%%%%%%%%%%%%%%%%%%%%%%%%%%%%%%%%%%%%%%%%%%%%%%%%%%%%%%%%%%%%%%%%%%%%%%%%%%%%%
%                         DATOS (Conectados a RMarkdown)                       %
%%%%%%%%%%%%%%%%%%%%%%%%%%%%%%%%%%%%%%%%%%%%%%%%%%%%%%%%%%%%%%%%%%%%%%%%%%%%%%%%
\title{Análisis Transaccional de venta minorista}
\author{Radamés Montiel \\ Rafael Ferrer}

% Lógica para Facultad: Si la defines en el YAML la usa, si no, usa la default

\facultad{Facultad de Ciencias Económicas y Sociales \\ Escuela de
Estadistica y Ciencias Actuariales}
  \escudofacultad{Latex/Classes/Escudos/faces}

\degree{Computación I}
\director{Prof.~Jesús Ochoa \\ Prof.~Oliver Triveño}
\degreedate{Febrero 2026} 
\lugar{Caracas}

\portadatrue 

% Metadatos PDF
\keywords{}
\subject{}

% --- CORRECCIÓN PARA PANDOC NUEVO (Imágenes) ---
\makeatletter
\providecommand{\pandocbounded}[1]{#1}
\makeatother


%%%%%%%%%%%%%%%%%%%%%%%%%%%%%%%%%%%%%%%%%%%%%%%%%%%%%
%                   DOCUMENTO                       %
%%%%%%%%%%%%%%%%%%%%%%%%%%%%%%%%%%%%%%%%%%%%%%%%%%%%%
\begin{document}

\maketitle

%%%%%%%%%%%%%%%%%%%%%%%%%%%%%%%%%%%%%%%%%%%%%%%%%%%%%
%                  PRÓLOGO                          %
%%%%%%%%%%%%%%%%%%%%%%%%%%%%%%%%%%%%%%%%%%%%%%%%%%%%%
\frontmatter

% Puedes agregar dedicatorias desde el YAML si quieres, o dejarlas fijas aquí

%%%%%%%%%%%%%%%%%%%%%%%%%%%%%%%%%%%%%%%%%%%%%%%%%%%%%
%                   ÍNDICES                         %
%%%%%%%%%%%%%%%%%%%%%%%%%%%%%%%%%%%%%%%%%%%%%%%%%%%%%
\setcounter{secnumdepth}{3}
\setcounter{tocdepth}{3}

\tableofcontents



\cleardoublepage

  \cleardoublepage       % Empieza en página derecha (estándar en tesis)
  \chapter*{Resumen}     % Crea el título "Resumen" sin numeración (Capítulo *)
  \addcontentsline{toc}{chapter}{Resumen} % (Opcional) Lo añade al índice
  
  En el presente informe se realiza un análisis exploratorio y
descriptivo de una base de datos transaccional del sector retail. Se
examina el impacto de las características demográficas de los clientes
(edad y género) sobre el comportamiento de compra, la preferencia de
categorías de productos y el volumen total de facturación. Mediante la
aplicación de técnicas de estadística descriptiva, segmentación de datos
y visualización interactiva, el estudio identifica los patrones de
consumo clave y perfila al cliente de mayor rentabilidad para respaldar
la toma de decisiones
comerciales.             % Aquí se pega el texto que escribiste en el YAML
  
  \cleardoublepage

%%%%%%%%%%%%%%%%%%%%%%%%%%%%%%%%%%%%%%%%%%%%%%%%%%%%%
%                   CONTENIDO (El Corazón)          %
%%%%%%%%%%%%%%%%%%%%%%%%%%%%%%%%%%%%%%%%%%%%%%%%%%%%%
\mainmatter
\def\baselinestretch{1.5}

% --- CAMBIO CRÍTICO 2: Aquí RMarkdown inyecta todo tu texto ---
\chapter*{Introducción}

Rentabilidad de la venta minorista mediante el análisis de la
segmentación demográfica y el comportamiento de las transacciones

En el sector minorista, conocer el volumen general de ventas no es
suficiente; el verdadero valor estratégico radica en entender qué
perfiles de clientes generan esas ventas.

En lo que a ventas se refiere, poder conocer no solo la cantidad de
ventas sino también qué clientes las produce, genera el mayor valor para
una empresa, ya que esta información puede llevar a que se tomen
decisiones en función de permitir la entrada de más clientes en los
sectores con mejor desempeño. Estas decisiones no solo están basadas en
la intuición, sino también basadas en datos, lo que podrá ser de gran
ayuda en la inversión en marketing y optimización del Retorno de
Inversión (ROI).

Por todo lo mencionado, el siguiente informe busca responder la
pregunta: ¿Cuáles son los grupos demográficos que generan mayor
rentabilidad y más transacciones en la tienda minorista?.

\section*{Estructura del Informe}

Este documento técnico se encuentra estructurado en cinco capítulos
fundamentales:

\begin{itemize}
    \item En el \textbf{Capítulo I}, se plantea el problema de investigación.
    
    \item El \textbf{Capítulo II} establece el Marco Teórico, definiendo los conceptos estadísticos clave como medidas de tendencia central, medidas de dispersión y correlación, necesarios para la interpretación de los datos.
    
    \item El \textbf{Capítulo III} describe el Marco Metodológico, detallando el proceso de limpieza de datos y los algoritmos utilizados para validar la calidad de la información.
    
    \item El \textbf{Capítulo IV} presenta el Análisis de Resultados, donde se exhiben las tablas y gráficos generados dinámicamente.
    
    \item Finalmente, en el \textbf{Capítulo V}, se exponen las Conclusiones y Recomendaciones estratégicas derivadas del análisis cuantitativo, orientadas a la toma de decisiones basada en datos (\textit{Data-Driven Decision Making}).
\end{itemize}

\chapter{Planteamiento del Problema}

En el entorno altamente competitivo del sector minorista (retail), el
éxito financiero de una empresa ya no depende únicamente de la cantidad
de productos que vende, sino de su capacidad para comprender a
profundidad quiénes son los clientes que generan dichas ventas.
Actualmente, las empresas recopilan un volumen considerable de datos
transaccionales diarios que incluyen perfiles demográficos (edad y
género), hábitos de compra (categorías y cantidades) y m´etricas de
impacto financiero (precio unitario y monto total). Sin embargo, el
problema radica en que, sin un análisis adecuado, esta información se
mantiene como datos crudos. La empresa se enfrenta a una incertidumbre
estratégica al no tener claramente identificados los patrones de consumo
de sus clientes. Al no conocer cómo las variables demográficas impactan
la rentabilidad, se corre el riesgo de invertir en campañas de marketing
ineficientes, mantener un inventario desalineado con las preferencias
reales de los compradores y perder oportunidades valiosas de ventas
cruzadas y fidelización de los segmentos de mayor valor. Ante esta
situación, surgen interrogantes fundamentales: ¿Son los jóvenes o los
adultos quienes sostienen financieramente el negocio? ¿Existe una brecha
de género en el ticket promedio de compra? La ausencia de respuestas
basadas en datos a estas interrogantes genera ineficiencias operativas y
financieras.

Argumento final. Debido a esto, se plantea las siguientes preguntas de
investigación:

\textbf{1.} ¿Varía el ticket promedio significativamente entre hombres y
mujeres?

\textbf{2.} ¿Cómo influye la edad de manera directa en la elección de la
categoría del producto?

\textbf{3.} Mientras mayor edad tienen los clientes, ¿hay tendencia a
comprar artículos premium y a un mayor gasto total?

\section{Justificación}

La justificación de este trabajo de investigación reside en la necesidad
de aplicar técnicas de análisis de datos e Inteligencia de Negocios para
transformar registros transaccionales en conocimiento estratégico. Desde
un punto de vista práctico, los resultados de este estudio permitirán a
la gerencia pasar de la intuición a la toma de decisiones basadas en
datos. Al identificar con exactitud el perfil del Cliente Ideal) y
comprender la rentabilidad segmentada, la empresa podrá optimizar su
Retorno de Inversión (ROI), focalizar sus presupuestos de marketing en
los segmentos correctos y adaptar el surtido de productos (económicos
vs.~premium).

\section{Objetivos}

\subsection{Objetivo General}

Analizar los patrones demográficos y el comportamiento de compra en los
registros transaccionales del sector retail, con el fin de caracterizar
su impacto en la facturación y proporcionar insights que permitan
optimizar las estrategias comerciales de la empresa.

\subsection{Objetivos Específicos}

\begin{enumerate}
  \item Determinar si existe una variación significativa en el ticket promedio de compra (Total Amount) entre hombres y mujeres.
  \item Evaluar de qué manera influye la edad y género del cliente en la elección de la categoría del producto (Ropa, Belleza, Electrónica).
  \item Contrastar la hipótesis de tendencia de compra Premium, verificando si a mayor edad los clientes tienden a comprar artículos de lujo (> \$300) y generar un mayor gasto total.
  \item Identificar qué rango de edad representa el mayor volumen de ingresos acumulados (motor financiero) para el negocio.
  \item Comprobar si el segmento demográfico que realiza la mayor cantidad de transacciones (volumen operativo) es el mismo que deja la mayor cantidad de dinero (volumen financiero).
\end{enumerate}

\section{Cobertura}

\subsection{Horizontal}

La cobertura horizontal abarca las características intrínsecas de cada
transacción registrada en la base de datos. Se consideran las
dimensiones demográficas del cliente (edad y género), las dimensiones
descriptivas del producto (categoría) y las métricas financieras
(cantidad comprada, precio por unidad y monto total facturado).

\begin{table}[ht]
\centering
\caption{Variables consideradas en el estudio} 
\resizebox{12cm}{!} {
\begin{tabular}{cll}
  \hline
  \textbf{TIPO} & \textbf{VARIABLE} & \textbf{DESCRIPCIÓN} \\ 
  \hline
Temporal & Date & Fecha exacta en la que se efectuó la transacción \\
Categórica & Gender & Género del cliente (Male, Female) \\
Numérica & Age & Edad biológica del comprador \\
Categórica & Product Category & Categoría del artículo comprado (e.g., Beauty, Clothing) \\
Numérica & Quantity & Cantidad de artículos llevados en una misma transacción \\
Numérica & Price per Unit & Precio unitario del producto \\
Numérica & Total Amount & Monto total facturado en la compra (Ticket) \\
   \hline
\end{tabular}
}
\end{table}

\subsection{Vertical}

La cobertura vertical del trabajo de investigación comprende un total de
1,000 observaciones (registros) válidos, tras el proceso de limpieza y
omisión de variables no analíticas (como Transaction ID y Customer ID).
Cada fila representa una compra única realizada por un cliente en la
tienda.

\section{Periodo de Referencia}

Los datos analizados corresponden a un historial transaccional de un año
calendario, abarcando las ventas registradas entre Enero de 2023 y
Diciembre de 2023. Esto permite capturar un ciclo comercial completo y
evitar sesgos de meses específicos.

\section{Variables}

La variable principal en estudio es \textbf{Total Amount}, que
representa el éxito financiero del negocio medido a través del ingreso
total generado por cada transacción.

Como variables explicativas o independientes se analizan:

\begin{enumerate}
    \item Variables Categóricas: \textbf{Gender} y \textbf{Product Category}
    \item Variables Numéricas: \textbf{Age}, \textbf{Quantity} y \textbf{Price per Unit}
    \item Variables Temporales: \textbf{Date}
\end{enumerate}

\chapter{Marco Teórico}

\section{Bases Teóricas}

El presente capítulo establece los fundamentos matemáticos y
estadísticos aplicados para el análisis de los datos transaccionales del
sector retail. Se aborda la teoría desde una perspectiva de inteligencia
de negocios, priorizando los conceptos que permiten perfilar al cliente,
medir el volumen financiero y calcular preferencias de consumo.

La comprensión de estos conceptos es el sustento metodológico de las
operaciones computacionales realizadas sobre el conjunto de datos para
dar respuesta a los objetivos de la investigación.

\subsection{Análisis Exploratorio de Datos (EDA)}

El Análisis Exploratorio de Datos es el paso preliminar para resumir las
características principales de un conjunto de transacciones.
Matemáticamente, para una variable numérica \(X\) (como el monto total
de la compra o \textit{Total Amount}), el análisis se fundamenta en sus
métricas agregadas. Sea \(x_1, x_2, \dots, x_n\) una muestra de ventas
observadas, definimos:

\begin{itemize}
    \item \textbf{Promedio:} Representa el ticket promedio esperado por cada transacción realizada.

\item \textbf{Mediana:} Representa el 50% central de toda la distribución. El 50% de las transacciones están por debajo de este valor, y el otro 50% están por encima. Es importante utilizar esta medida ya que nos ayuda a saber el cómo están distribuidos los ingresos alrededor del valor central.

\item \textbf{Desviación estándar:} Evalúa qué tan dispersos están los datos alrededor de la media, si estos están muy cerca, significa que no hay mucha dispersión, si están muy lejos, significa que hay mucha dispersión. Es importante utilizar esta medida ya que nos permite saber si el gasto promedio por algún grupo está muy disperso o poco disperso.

\item \textbf{Cuartiles y Box Plot:} Los cuartiles separan la muestra ordenada en cuatro partes iguales. El Box Plot grafica estos cuartiles permitiendo saber qué tanta asimetría hay en una distribución, también permite evaluar el 50% central (Rango Inter-Cuartilico) e identificar valores atípicos. Es de suma importancia utilizar esta medida ya que nos ayudará a saber sobre la variabilidad y forma del gasto de cualquier grupo.

\item \textbf{Correlación Lineal:} Nos permite visualizar si existe alguna correlación entre 2 variables numéricas, si la correlación es positiva, significa que los valores se mueven en la misma dirección, si es negativa, se mueven en direcciones opuestas. Es de vital importancia utilizar esta medida ya que nos permite saber si, por ejemplo, el gasto y la frecuencia de compra están correlacionadas.

\item \textbf{Discretización de variables continuas}: Como el mercado minorista no es homogéneo, por lo que para poder analizar este sectos se tuvo que separar por rangos y segmentos las variables a estudiar, por ejemplo:

En el caso de "Precio", se separó por compras Económicas y por compras Premium.

En el caso de "Edad", se separó por cinco Segmentos Etarios.

Todo esto para facilitar el análisis y encontrar patrones en los datos.

\end{itemize}

\chapter{Marco Metódico}

En esta sección se presentan las metodologías, procedimientos
algorítmicos y técnicas descriptivas implementadas para dar respuesta a
los objetivos planteados en el capítulo anterior. Se detalla el
tratamiento de la fuente de datos transaccional, las técnicas de
limpieza (Data Wrangling), la transformación de variables y los métodos
de agrupación utilizados para caracterizar el comportamiento de los
clientes.

El análisis se ha realizado utilizando el lenguaje de programación
\textbf{R} (versión 4.x) bajo el entorno de desarrollo RStudio,
empleando el paradigma \textbf{Tidyverse}, específicamente, la librería
dplyr para la manipulación de datos estructurados, lubridate para la
manipuación de fechas, y ggplot2 para la visualización de los
resultados.

\section{Bases metódicas utilizadas en el trabajo de investigación}

El diseño de la investigación es de carácter
\textbf{cuantitativo, no experimental y descriptivo-exploratorio}, dado
que se analizan datos históricos observados sin manipulación directa de
las variables. A continuación, se describen las etapas del procesamiento
analítico.

\subsection{Procesamiento y Limpieza de Datos}

La calidad de los resultados depende directamente de la integridad de
los datos de entrada. El proceso de limpieza consistió en la inspección
de valores nulos (\textit{missing values}), detección de registros
duplicados, outliers y la adaptación de los tipos de datos.

Específicamente, se verificó mediante funciones de R que la base de
datos original se encontraba en perfectas condiciones operativas. Para
optimizar el procesamiento en memoria, se procedió a la eliminación de
las columnas identificadoras (\texttt{Customer ID} y
\texttt{Transaction ID}) que carecen de valor analítico, se estandarizó
la variable temporal al formato de fecha correcto (\texttt{Date}) y se
tradujeron los nombres de las variables al español

\begin{Shaded}
\begin{Highlighting}[]
\FunctionTok{library}\NormalTok{(tidyverse)}
\FunctionTok{library}\NormalTok{(knitr)}

\NormalTok{datos\_retail }\OtherTok{\textless{}{-}} \FunctionTok{read\_csv}\NormalTok{(}\StringTok{"venta{-}minorista.csv"}\NormalTok{)}

\CommentTok{\# Buscando datos nulos y duplicados}

\FunctionTok{sum}\NormalTok{(}\FunctionTok{is.na}\NormalTok{(datos\_retail))}

\FunctionTok{anyDuplicated}\NormalTok{(datos\_retail)}

\CommentTok{\# Clases de las columnas}

\FunctionTok{str}\NormalTok{(datos\_retail)}

\CommentTok{\# Verificando outliers}

\FunctionTok{table}\NormalTok{(datos\_retail}\SpecialCharTok{$}\NormalTok{Quantity}\SpecialCharTok{*}\NormalTok{datos\_retail}\SpecialCharTok{$}\StringTok{\textasciigrave{}}\AttributeTok{Price per Unit}\StringTok{\textasciigrave{}} \SpecialCharTok{==}\NormalTok{ datos\_retail}\SpecialCharTok{$}\StringTok{\textasciigrave{}}\AttributeTok{Total Amount}\StringTok{\textasciigrave{}}\NormalTok{)}

\FunctionTok{boxplot}\NormalTok{(datos\_retail}\SpecialCharTok{$}\NormalTok{Age)}
\end{Highlighting}
\end{Shaded}

\begin{Shaded}
\begin{Highlighting}[]
\FunctionTok{boxplot.stats}\NormalTok{(datos\_retail}\SpecialCharTok{$}\NormalTok{Age)}
\FunctionTok{boxplot}\NormalTok{(datos\_retail}\SpecialCharTok{$}\NormalTok{Quantity)}
\end{Highlighting}
\end{Shaded}

\begin{Shaded}
\begin{Highlighting}[]
\FunctionTok{boxplot.stats}\NormalTok{(datos\_retail}\SpecialCharTok{$}\NormalTok{Quantity)}
\FunctionTok{boxplot}\NormalTok{(datos\_retail}\SpecialCharTok{$}\StringTok{\textasciigrave{}}\AttributeTok{Price per Unit}\StringTok{\textasciigrave{}}\NormalTok{)}
\end{Highlighting}
\end{Shaded}

\begin{Shaded}
\begin{Highlighting}[]
\FunctionTok{boxplot.stats}\NormalTok{(datos\_retail}\SpecialCharTok{$}\StringTok{\textasciigrave{}}\AttributeTok{Price per Unit}\StringTok{\textasciigrave{}}\NormalTok{)}
\FunctionTok{boxplot}\NormalTok{(datos\_retail}\SpecialCharTok{$}\StringTok{\textasciigrave{}}\AttributeTok{Total Amount}\StringTok{\textasciigrave{}}\NormalTok{)}
\end{Highlighting}
\end{Shaded}

\begin{Shaded}
\begin{Highlighting}[]
\FunctionTok{boxplot.stats}\NormalTok{(datos\_retail}\SpecialCharTok{$}\StringTok{\textasciigrave{}}\AttributeTok{Total Amount}\StringTok{\textasciigrave{}}\NormalTok{)}

\FunctionTok{table}\NormalTok{(datos\_retail}\SpecialCharTok{$}\NormalTok{Gender)}
\FunctionTok{table}\NormalTok{(datos\_retail}\SpecialCharTok{$}\StringTok{\textasciigrave{}}\AttributeTok{Product Category}\StringTok{\textasciigrave{}}\NormalTok{)}

\CommentTok{\# Convirtiendo la columna Date de texto a formato fecha}

\FunctionTok{library}\NormalTok{(lubridate)}
\NormalTok{datos\_retail}\SpecialCharTok{$}\NormalTok{Date }\OtherTok{\textless{}{-}} \FunctionTok{ymd}\NormalTok{(datos\_retail}\SpecialCharTok{$}\NormalTok{Date)}
\FunctionTok{class}\NormalTok{(datos\_retail}\SpecialCharTok{$}\NormalTok{Date)}

\CommentTok{\# Cambiando el nombre de las variables a español}

\NormalTok{datos\_retail }\OtherTok{\textless{}{-}}\NormalTok{ datos\_retail }\SpecialCharTok{\%\textgreater{}\%}
  \FunctionTok{rename}\NormalTok{(}
    \AttributeTok{ID\_Transaccion =} \StringTok{\textasciigrave{}}\AttributeTok{Transaction ID}\StringTok{\textasciigrave{}}\NormalTok{,}
    \AttributeTok{Fecha =}\NormalTok{ Date,}
    \AttributeTok{Cliente\_ID =} \StringTok{\textasciigrave{}}\AttributeTok{Customer ID}\StringTok{\textasciigrave{}}\NormalTok{,}
    \AttributeTok{Genero =}\NormalTok{ Gender,}
    \AttributeTok{Edad =}\NormalTok{ Age,}
    \AttributeTok{Categoria\_producto =} \StringTok{\textasciigrave{}}\AttributeTok{Product Category}\StringTok{\textasciigrave{}}\NormalTok{,}
    \AttributeTok{Cantidad =}\NormalTok{ Quantity,}
    \AttributeTok{Precio\_por\_cantidad =} \StringTok{\textasciigrave{}}\AttributeTok{Price per Unit}\StringTok{\textasciigrave{}}\NormalTok{,}
    \AttributeTok{Monto\_total =} \StringTok{\textasciigrave{}}\AttributeTok{Total Amount}\StringTok{\textasciigrave{}}
\NormalTok{  )}

\NormalTok{datos\_retail }\OtherTok{\textless{}{-}}\NormalTok{ datos\_retail }\SpecialCharTok{\%\textgreater{}\%}
  \FunctionTok{mutate}\NormalTok{(}\AttributeTok{Genero =} \FunctionTok{case\_match}\NormalTok{(Genero,}
                           \StringTok{"Female"} \SpecialCharTok{\textasciitilde{}} \StringTok{"Mujer"}\NormalTok{,}
                         \StringTok{"Male"} \SpecialCharTok{\textasciitilde{}} \StringTok{"Hombre"}\NormalTok{),}
         \AttributeTok{Categoria\_producto =} \FunctionTok{case\_match}\NormalTok{(Categoria\_producto,}
                                         \StringTok{"Clothing"} \SpecialCharTok{\textasciitilde{}} \StringTok{"Ropa"}\NormalTok{,}
                                         \StringTok{"Electronics"} \SpecialCharTok{\textasciitilde{}} \StringTok{"Electronica"}\NormalTok{,}
                                         \StringTok{"Beauty"} \SpecialCharTok{\textasciitilde{}} \StringTok{"Belleza"}\NormalTok{))}

\NormalTok{datos\_limpios }\OtherTok{\textless{}{-}}\NormalTok{ datos\_retail }\SpecialCharTok{\%\textgreater{}\%}
  \FunctionTok{drop\_na}\NormalTok{() }\SpecialCharTok{\%\textgreater{}\%}
  \FunctionTok{distinct}\NormalTok{()}

\FunctionTok{write.csv}\NormalTok{(datos\_limpios, }\StringTok{"retail\_limpio.csv"}\NormalTok{, }\AttributeTok{row.names =} \ConstantTok{FALSE}\NormalTok{)}
\end{Highlighting}
\end{Shaded}

\subsection{Transformación y Discretización de Variables (Feature Engineering)}

Para habilitar el análisis segmentado (estratificación) requerido por
los objetivos de la investigación, fue estrictamente necesario
transformar variables numéricas continuas en variables categóricas
ordinales. Este proceso de discretización se realizó empleando la
función \texttt{mutate()} combinada con sentencias lógicas condicionales
(\texttt{ifelse()}).

Se crearon dos nuevas dimensiones analíticas fundamentales en el
conjunto de datos:

\begin{itemize}
\item \textbf{Grupo Etario:} Para contrastar el comportamiento generacional, se agruparon a los clientes en un rango de edad de aproximadamente de 8 a 10 años de diferencia por segmento etario: 

Adulto joven (18-26): Entrada al mercado laboral o estudiantes.

Adulto (27-35): Representan a los jóvenes profesionales con una mayor estabilidad económica que la del grupo anterior.

Adulto maduro (36-43): Representa un periodo de mayor establidad económica que el grupo anterior, además de que poseen mayor experiencia laboral.

Mediana edad (44-53): Representa el periodo de consolidación profesional y económica.

Adulto mayor (54-44): Representan a los consumidores que entran en la etapa del retiro laboral.

\item \textbf{Tipo de Producto:} Para clasificar los artículos según su precio por unidad en gamas "Económicas" ($\leq \$300$) y "Premium" ($> \$300$).

Económicas: Representan un mayor volumen operativo en el negocio. Son compras que gastan menos pero tienen una mayor frecuencia que las compras Premium.

Premium: Representan de los mayores ingresos del negocio. Son compras que a pesar de que son menos frecuentes, estas compras debido a su alto valor generan una cantidad de ingresos importante para el negocio.

\end{itemize}

\begin{Shaded}
\begin{Highlighting}[]
\CommentTok{\# Creación de variables categóricas condicionales}
\NormalTok{retail }\OtherTok{\textless{}{-}} \FunctionTok{read\_csv}\NormalTok{(}\StringTok{"retail\_limpio.csv"}\NormalTok{)}

\FunctionTok{max}\NormalTok{(retail}\SpecialCharTok{$}\NormalTok{Edad)}
\FunctionTok{min}\NormalTok{(retail}\SpecialCharTok{$}\NormalTok{Edad)}

\NormalTok{cortes }\OtherTok{\textless{}{-}} \FunctionTok{c}\NormalTok{(}\DecValTok{18}\NormalTok{, }\DecValTok{27}\NormalTok{, }\DecValTok{36}\NormalTok{, }\DecValTok{45}\NormalTok{, }\DecValTok{54}\NormalTok{, }\DecValTok{64}\NormalTok{)}
\NormalTok{etiquetas }\OtherTok{\textless{}{-}} \FunctionTok{c}\NormalTok{(}\StringTok{"Adulto joven (18{-}26)"}\NormalTok{,}
               \StringTok{"Adulto (27{-}35)"}\NormalTok{,}
               \StringTok{"Adulto Maduro (36{-}43)"}\NormalTok{,}
               \StringTok{"Mediana edad (44{-}53)"}\NormalTok{,}
               \StringTok{"Adulto Mayor (54{-}64)"}\NormalTok{)}

\NormalTok{retail[}\StringTok{"Rango\_edad"}\NormalTok{] }\OtherTok{\textless{}{-}} \FunctionTok{cut}\NormalTok{(retail}\SpecialCharTok{$}\NormalTok{Edad,}
                   \AttributeTok{breaks =}\NormalTok{ cortes,}
                   \AttributeTok{labels =}\NormalTok{ etiquetas,}
                   \AttributeTok{include.lowest =} \ConstantTok{TRUE}\NormalTok{)}
\CommentTok{\# Acumulado de gasto por mes}

\NormalTok{datos\_mensual }\OtherTok{\textless{}{-}}\NormalTok{ retail }\SpecialCharTok{\%\textgreater{}\%}
  \FunctionTok{mutate}\NormalTok{(}\AttributeTok{mes =} \FunctionTok{floor\_date}\NormalTok{(Fecha, }\StringTok{"month"}\NormalTok{)) }\SpecialCharTok{\%\textgreater{}\%}
  \FunctionTok{group\_by}\NormalTok{(mes, Rango\_edad) }\SpecialCharTok{\%\textgreater{}\%}
  \FunctionTok{summarise}\NormalTok{(}\AttributeTok{gasto\_por\_mes =} \FunctionTok{sum}\NormalTok{(Monto\_total), }\AttributeTok{.groups =} \StringTok{"drop"}\NormalTok{)}
 
\CommentTok{\# Gasto acumulado al pasar de los meses}

\NormalTok{datos\_mensual }\OtherTok{\textless{}{-}}\NormalTok{ datos\_mensual }\SpecialCharTok{\%\textgreater{}\%}
  \FunctionTok{arrange}\NormalTok{(mes) }\SpecialCharTok{\%\textgreater{}\%}
  \FunctionTok{group\_by}\NormalTok{(Rango\_edad) }\SpecialCharTok{\%\textgreater{}\%}
  \FunctionTok{mutate}\NormalTok{(}\AttributeTok{Gasto\_acumulado =} \FunctionTok{cumsum}\NormalTok{(gasto\_por\_mes)) }\SpecialCharTok{\%\textgreater{}\%}
  \FunctionTok{ungroup}\NormalTok{() }\SpecialCharTok{\%\textgreater{}\%}
  \FunctionTok{filter}\NormalTok{(mes }\SpecialCharTok{!=} \StringTok{"2024{-}01{-}01"}\NormalTok{)}

\CommentTok{\# Para saber el porcentaje por la compra de categoria de edad por segmento etario}

\NormalTok{preferencias }\OtherTok{\textless{}{-}}\NormalTok{ retail }\SpecialCharTok{\%\textgreater{}\%}
  \FunctionTok{group\_by}\NormalTok{(Genero, Rango\_edad, Categoria\_producto) }\SpecialCharTok{\%\textgreater{}\%}
  \FunctionTok{summarise}\NormalTok{(}\AttributeTok{total\_elecciones =} \FunctionTok{n}\NormalTok{(), }\AttributeTok{.groups =} \StringTok{"drop"}\NormalTok{) }\SpecialCharTok{\%\textgreater{}\%}
  \FunctionTok{mutate}\NormalTok{(}\AttributeTok{Porcentaje\_preferencia =}\NormalTok{ (total\_elecciones }\SpecialCharTok{/} \FunctionTok{sum}\NormalTok{(total\_elecciones) }\SpecialCharTok{*} \DecValTok{100}\NormalTok{)) }\SpecialCharTok{\%\textgreater{}\%} 
  \FunctionTok{group\_by}\NormalTok{(Rango\_edad) }\SpecialCharTok{\%\textgreater{}\%}
  \FunctionTok{mutate}\NormalTok{(}\AttributeTok{Porcentaje\_por\_segmento =}\NormalTok{ (total\_elecciones }\SpecialCharTok{/} \FunctionTok{sum}\NormalTok{(total\_elecciones)) }\SpecialCharTok{*} \DecValTok{100}\NormalTok{) }\SpecialCharTok{\%\textgreater{}\%}
  \FunctionTok{ungroup}\NormalTok{()}

\NormalTok{retail }\OtherTok{\textless{}{-}}\NormalTok{ retail }\SpecialCharTok{\%\textgreater{}\%}
  \FunctionTok{mutate}\NormalTok{(}\AttributeTok{segmento =} \FunctionTok{ifelse}\NormalTok{(Precio\_por\_cantidad }\SpecialCharTok{\textgreater{}=} \DecValTok{300}\NormalTok{, }\StringTok{"Premium"}\NormalTok{, }\StringTok{"Economico"}\NormalTok{))}

\NormalTok{perfiles\_premium }\OtherTok{\textless{}{-}}\NormalTok{ retail }\SpecialCharTok{\%\textgreater{}\%}
  \FunctionTok{group\_by}\NormalTok{(Rango\_edad, segmento) }\SpecialCharTok{\%\textgreater{}\%}
  \FunctionTok{summarise}\NormalTok{(}
    \AttributeTok{Total\_ventas =} \FunctionTok{n}\NormalTok{(),}
    \AttributeTok{Monto\_total =} \FunctionTok{sum}\NormalTok{(Monto\_total),}
    \AttributeTok{.groups =} \StringTok{\textquotesingle{}drop\textquotesingle{}}
\NormalTok{  ) }\SpecialCharTok{\%\textgreater{}\%}
  \FunctionTok{group\_by}\NormalTok{(Rango\_edad) }\SpecialCharTok{\%\textgreater{}\%}
  \FunctionTok{mutate}\NormalTok{(}\AttributeTok{Porcentaje\_por\_segmento =}\NormalTok{ (Monto\_total }\SpecialCharTok{/} \FunctionTok{sum}\NormalTok{(Monto\_total)) }\SpecialCharTok{*} \DecValTok{100}\NormalTok{) }\SpecialCharTok{\%\textgreater{}\%}
  \FunctionTok{ungroup}\NormalTok{() }\SpecialCharTok{\%\textgreater{}\%}
  \FunctionTok{mutate}\NormalTok{(}\AttributeTok{Porcentaje\_global =}\NormalTok{ (Monto\_total }\SpecialCharTok{/} \FunctionTok{sum}\NormalTok{(Monto\_total)) }\SpecialCharTok{*} \DecValTok{100}\NormalTok{)}
  
\NormalTok{resumen\_gastos }\OtherTok{\textless{}{-}}\NormalTok{ retail }\SpecialCharTok{\%\textgreater{}\%}
  \FunctionTok{group\_by}\NormalTok{(Rango\_edad) }\SpecialCharTok{\%\textgreater{}\%}
  \FunctionTok{summarise}\NormalTok{(}\AttributeTok{Gasto\_total\_grupo =} \FunctionTok{sum}\NormalTok{(Monto\_total)) }
\end{Highlighting}
\end{Shaded}

\chapter{Análisis de Resultados}

\small Este capítulo presenta los resultados obtenidos al aplicar las
metodologías descritas en el capítulo anterior. Se expone el análisis
estadístico descriptivo de las variables estudiadas, validando las
hipótesis de negocio formuladas respecto al comportamiento demográfico y
financiero de los clientes.

\section{Ticket Promedio por Género}

El primer objetivo del análisis consistió en determinar si existía una
variación significativa en el ticket promedio de compra
(\textit{Total Amount}) entre hombres y mujeres, lo cual es vital para
la asignación de presupuestos de marketing.

Como se evidencia en el Cuadro \ref{tab:tabla_genero}, el ticket
promedio no varía significativamente entre géneros (Mujeres: \$457 vs
Hombres: \$455). La diferencia es mínima, lo que indica que el género no
es un factor determinante en la disposición de gasto por visita.

\begin{table}

\caption{\label{tab:tabla_genero}Estadísticas descriptivas del monto total por segmento etario}
\centering
\begin{tabular}[t]{lrrrr}
\toprule
Género & Número de compras & Promedio & Mediana & Desviación estándar\\
\midrule
Hombre & 490 & 455 & 120 & 566\\
Mujer & 510 & 457 & 150 & 554\\
\bottomrule
\end{tabular}
\end{table}

\begin{figure}[H]

{\centering \includegraphics[width=1\linewidth,]{informe_files/figure-latex/obj1_plot_barras-1} 

}

\caption{Ticket Promedio de Compra por Género}\label{fig:obj1_plot_barras}
\end{figure}

Como se puede observar en la figura \ref{fig:obj1_plot_barras} que
hombres y mujeres gastan una cantidad de dinero similar. Los hombres
tienen una mayor predilección por el apartado de ``Electrónica'',
mientras que las mujeres se decantan más por el apartado de ``Ropa''.
Aunque esto no significa que se deban descuidar los otros productos de
ambos géneros por prejuicios, por ejemplo, la barra más baja de toda la
distribución es la de los hombres en el apartado de ``Belleza'', pero
este apartado roza los \$70,000, casi igualando la cantidad gastadas por
mujeres en ese mismo apartado, por lo que no se deben descuidar las
ventas en ningún nicho de la distribución.

\begin{figure}[H]

{\centering \includegraphics[width=1\linewidth,]{informe_files/figure-latex/obj1_plot_boxplot-1} 

}

\caption{Ticket Promedio de Compra por Género}\label{fig:obj1_plot_boxplot}
\end{figure}

La cantidad mediana de gasto en ambos géneros no varía mucho, siendo que
la mediana del gasto de los hombres es de \$120 y la de las mujeres es
de \$150. Al ser la caja del mismo tamaño en ambos casos, esto sugiere
que ambos géneros tienen una variabilidad parecida, lo que indica que
son igual de predecibles en sus compras. Se puede notar como el precio
medio está muy comprimido en la parte inferior de la caja, por lo que el
50\% que menos gastan están muy comprimidos en la parte inferior de la
caja en un pequeño rango de \$0 a \$120.

Se puede observar como el 75\% de los clientes gastan menos de 900, por
lo que variabilidad es predecible en ambos casos, esto significa también
que el techo son los 900, por lo que pasar de esta barrera va a ser un
poco más difícil ya que pasa de lo que es el precio común de la compra.

\begin{figure}[H]

{\centering \includegraphics[width=1\linewidth,]{informe_files/figure-latex/obj1_plot_densidad-1} 

}

\caption{Ticket Promedio de Compra por Género}\label{fig:obj1_plot_densidad}
\end{figure}

En el gráfico de la figura \ref{fig:obj1_plot_densidad} se puede
observar la forma de la distribución del gasto de hombres y mujeres. Al
estar la curva sesgada hacia la derecha, lo que quiere decir que el
promedio de compra es mayor que las compras que más se realizan, esto
significa que la gran mayoría de las compras entre hombres y mujeres
están situadas en el ticket más bajo. Se puede observar un
comportamiento homogéneo de las compras entre hombres y mujeres, esto
indica que el género no influye tanto en la compra de productos, por lo
que se debería enfocar más las ventas hacia los productos que compran
ambos géneros, en vez de qué género compra.

En el final de la curva, hay un pequeño repunte en los \$2000, estos
aunque sean pocos en la distribución, representan un gran valor debido a
la cantidad que gastan, por lo que se puede identificar la clase de
productos que compran y darles un trato preferencial a partir de estos
productos.

\clearpage

\section{Análisis del gasto acumulado por Segmento Etario}

.

\begin{figure}[H]

{\centering \includegraphics[width=1\linewidth,]{informe_files/figure-latex/obj2_plot_calor-1} 

}

\caption{Preferencias de Categoría según Grupo Etario}\label{fig:obj2_plot_calor}
\end{figure}

El análisis de la figura \ref{fig:obj2_plot_calor} nos muestra patrones
para identificar las preferencias por segmento demográfico. Se puede
observar como las mujeres jóvenes tienen una mayor predilección a
comprar productos de Belleza (4.10\%), lo que significa que hay un nicho
de mercado altamente especializado en esta categoria.

Si observamos la fila de adultos mayores, se evidencia como las mujeres
mayores que compran ropa tienen el mayor porcentaje de toda la
distribución (4.20\%), a su vez, los hombres mayores tienen una mayor
inclinación a comprar productos de Electrónica (4.10\%). Esto indica
que, si bien ambos géneros tienen una cantidad de gasto similar, la
preferencia del consumidor está fuertemente ligada a la edad y el género
del cliente.

\begin{figure}[H]

{\centering \includegraphics[width=1\linewidth,]{informe_files/figure-latex/obj2_plot_calor2-1} 

}

\caption{Preferencias de Categoría según Grupo Etario}\label{fig:obj2_plot_calor2}
\end{figure}

En la figura \ref{fig:obj2_plot_calor2} se puede ver la intensidad de
consumo por segmento separado por su segmento etario. Se puede ver como
los hombres en cada una de sus distribuciones por segmento etario tienen
muy poco interés por los productos de belleza, siendo que todas sus
compras están por debajo de la mediana. En contraste, las mujeres tienen
una mayor predilección por los productos de belleza, aunque las mujeres
de 36-43 años tienen menor interés por estos productos a diferencia de
los segmentos etarios.

Otro hallazgo, es que las mujeres mayores de 36 años tienen una mayor
inclinación a comprar ropa que las mujeres menores de 36 años. También
se puede observar como los hombres adultos tienen una mayor predilección
a comprar ropa que cualquier otro producto, lo que significa que la ropa
representa un ``producto ancla'' para este segmento demográfico.

Se puede observar como mientras pasa la edad los hombres van tomando
mayor preferencia hacia comprar productos de Electrónica, siendo que los
hombres mediana edad que compraron productos electrónicos representan el
19.71\% de su distribución.

\clearpage

\section{Tendencia de Compra: Económico vs Premium}

Se contrastó la hipótesis general de que, a mayor edad, los clientes
tienen una mayor tendencia a adquirir artículos ``Premium'' (precio
unitario \textgreater{} \$300).

\begin{table}[!h]
\centering
\caption{\label{tab:tabla_premium}Estadísticas descriptivas de las compras Premium vs. Estándar por segmento etario}
\centering
\resizebox{\ifdim\width>\linewidth\linewidth\else\width\fi}{!}{
\begin{tabular}[t]{llrrrr}
\toprule
Segmento Etario & Tipo de compra & Número de compras & Promedio & Mediana & DE\\
\midrule
\cellcolor{gray!10}{Adulto joven (18-26)} & \cellcolor{gray!10}{Economico} & \cellcolor{gray!10}{121} & \cellcolor{gray!10}{80} & \cellcolor{gray!10}{60} & \cellcolor{gray!10}{51}\\
Adulto joven (18-26) & Premium & 93 & 1056 & 1000 & 520\\
\cellcolor{gray!10}{Adulto (27-35)} & \cellcolor{gray!10}{Economico} & \cellcolor{gray!10}{106} & \cellcolor{gray!10}{98} & \cellcolor{gray!10}{100} & \cellcolor{gray!10}{51}\\
Adulto (27-35) & Premium & 69 & 1070 & 1000 & 550\\
\cellcolor{gray!10}{Adulto Maduro (36-43)} & \cellcolor{gray!10}{Economico} & \cellcolor{gray!10}{118} & \cellcolor{gray!10}{83} & \cellcolor{gray!10}{68} & \cellcolor{gray!10}{53}\\
\addlinespace
Adulto Maduro (36-43) & Premium & 69 & 1058 & 900 & 517\\
\cellcolor{gray!10}{Mediana edad (44-53)} & \cellcolor{gray!10}{Economico} & \cellcolor{gray!10}{125} & \cellcolor{gray!10}{97} & \cellcolor{gray!10}{100} & \cellcolor{gray!10}{52}\\
Mediana edad (44-53) & Premium & 83 & 949 & 900 & 488\\
\cellcolor{gray!10}{Adulto Mayor (54-64)} & \cellcolor{gray!10}{Economico} & \cellcolor{gray!10}{134} & \cellcolor{gray!10}{83} & \cellcolor{gray!10}{75} & \cellcolor{gray!10}{48}\\
Adulto Mayor (54-64) & Premium & 82 & 965 & 950 & 507\\
\bottomrule
\end{tabular}}
\end{table}

\begin{figure}[H]

{\centering \includegraphics[width=1\linewidth,]{informe_files/figure-latex/obj3_plot-1} 

}

\caption{Proporción de compras Premium vs Económico por Edad}\label{fig:obj3_plot}
\end{figure}

Se puede observar en la figura \ref{fig:obj3_plot} un gráfico de barras
agrupadas para visualizar la compras Premium y Estándar. Analizando el
gráfico, se puede notar como las compras premium generan una cantidad de
ganancia por compras Premium es mucho mayor a la de compras Económicas.
Esto indica que aunque en la tabla de la figura \ref{tab:tabla_premium}
se muestra como la gran mayoría de compras son realizadas por menos de
\$300 (Económico), la mayor fuente de ingreso del negocio es la compra
de artículos Premium.

Esto refuta la hipótesis inicial, debido a que los adultos mayores no
son el grupo que más gasta en compras Premium, sino que al contrario,
los adultos jóvenes son los que más gastan en compras Premium y son de
los que menos gastan en compras Económicas.

\clearpage

\section{Ingresos Acumulados por Segmento Etario}

\begin{table}[!h]
\centering
\caption{\label{tab:tabla_resumen_edades}Estadísticas descriptivas del monto total por segmento etario}
\centering
\begin{tabular}[t]{lrrrr}
\toprule
Segmento Etario & Número de compras & Promedio & Mediana & DE\\
\midrule
Adulto joven (18-26) & 214 & 504 & 150 & 594\\
Adulto (27-35) & 175 & 481 & 150 & 589\\
Adulto Maduro (36-43) & 187 & 443 & 120 & 568\\
Mediana edad (44-53) & 208 & 437 & 150 & 521\\
Adulto Mayor (54-64) & 216 & 418 & 120 & 531\\
\bottomrule
\end{tabular}
\end{table}

\begin{figure}[H]

{\centering \includegraphics[width=1\linewidth,]{informe_files/figure-latex/obj4_plot_barras-1} 

}

\caption{Ingresos Acumulados por Segmento Etario}\label{fig:obj4_plot_barras}
\end{figure}

En el gráfico de la figura \ref{fig:obj4_plot_barras} se puede observar
la distribución de gasto de los distintos segmentos etarios. En el los
adultos jóvenes gastan una mayor cantidad de dinero. En la tabla de la
figura \ref{tab:tabla_resumen_edades}, como aunque los adultos de 27 a
35 años son el segmento etario que menos transacciones generan, son el
segundo segmento etario con la mayor media de gasto de de toda la
distribución.

\begin{figure}[H]

{\centering \includegraphics[width=1\linewidth,]{informe_files/figure-latex/obj4_plot_lineas-1} 

}

\caption{Ingresos Acumulados por Segmento Etario}\label{fig:obj4_plot_lineas}
\end{figure}

El gráfico de la figura \ref{fig:obj4_plot_lineas}, es un gráfico de
líneas acumuladas el cual muestra cuál mes fue en el que un Segmento
Etario tuvo mayor gasto. Este gráfico muestra como los adultos jóvenes
tienen el mayor gasto acumulado a finales del 2023 (\$100K), se puede
observar como estos tienen un mayor ritmo de consumo entre Enero-Febrero
y Octubre-Noviembre, donde se puede ver que la pendiente de la curva es
más pronunciada, lo que indica un mayor consumo en estos meses.

Por otro lado, se puede notar como los adultos mayores a 44 años tienen
una tendencia de gasto similar, terminando ambos el año con \$90K,
también su pendiente es muy parecida en el periodo Abril-Septiembre,
esto indica que ambos grupos tienen unos hábitos de gasto similares a
mediados de año lo que puede dar paso a una campaña de marketing a estos
Segmentos Etarios a mediados de año.

\begin{figure}[H]

{\centering \includegraphics[width=1\linewidth,]{informe_files/figure-latex/obj4_plot_boxplot-1} 

}

\caption{Ingresos Acumulados por Segmento Etario}\label{fig:obj4_plot_boxplot}
\end{figure}

Se puede observar en la figura \ref{fig:obj4_plot_boxplot} la
variabilidad por cada Segmento Etario. Analizando el gráfico, se puede
notar como todos los Segmentos etarios tienen una variabilidad similar,
siendo que la mediana para cada grupo es parecida y la media es mucho
mayor a la mediana, esto indica que la distribución tiene un fuerte
sesgo a la derecha, lo que a su vez significa que la mayoría de las
compras de todos los grupos se concentran en valores pequeños,
específicamente entre \$120 y \$150. Por otro lado, los adultos mayores
una caja mucho más compacta que, lo que sugiere una menor variabilidad
que los otros Segmentos Etarios. sin embargo, este grupo tiene dos
valores atípicos claros entre \$1500 y \$2000. En otros grupos ambos
valores aún se consideran dentro de lo esperado, pero específicamente en
este Segmento Etario, este tipo de compras son excepcionales y que se
alejan mucho de su patrón estándar.

\clearpage

\section{Correlación entre Volumen Operativo y Financiero}

Finalmente, se cruzaron las variables de Rango de Edad y Género para
identificar de manera exacta el \textit{Buyer Persona} principal de la
tienda, evaluando si el segmento que realiza más transacciones es el que
genera mayores ingresos.

\begin{table}[!h]
\centering
\caption{\label{tab:tabla_segmento_demografico}Estadísticas descriptivas del monto total por segmento etario}
\centering
\resizebox{\ifdim\width>\linewidth\linewidth\else\width\fi}{!}{
\begin{tabular}[t]{llrrrrr}
\toprule
Género & Segmento Etario & Número de compras & Gasto total & Promedio & Mediana & Desviación estándar\\
\midrule
\cellcolor{gray!10}{Hombre} & \cellcolor{gray!10}{Adulto joven (18-26)} & \cellcolor{gray!10}{105} & \cellcolor{gray!10}{53790} & \cellcolor{gray!10}{512} & \cellcolor{gray!10}{150} & \cellcolor{gray!10}{597}\\
Hombre & Adulto (27-35) & 88 & 40680 & 462 & 150 & 597\\
\cellcolor{gray!10}{Hombre} & \cellcolor{gray!10}{Adulto Maduro (36-43)} & \cellcolor{gray!10}{88} & \cellcolor{gray!10}{39110} & \cellcolor{gray!10}{444} & \cellcolor{gray!10}{120} & \cellcolor{gray!10}{559}\\
Hombre & Mediana edad (44-53) & 102 & 44670 & 438 & 150 & 548\\
\cellcolor{gray!10}{Hombre} & \cellcolor{gray!10}{Adulto Mayor (54-64)} & \cellcolor{gray!10}{107} & \cellcolor{gray!10}{44910} & \cellcolor{gray!10}{420} & \cellcolor{gray!10}{100} & \cellcolor{gray!10}{539}\\
\addlinespace
Mujer & Adulto joven (18-26) & 109 & 54125 & 497 & 150 & 594\\
\cellcolor{gray!10}{Mujer} & \cellcolor{gray!10}{Adulto (27-35)} & \cellcolor{gray!10}{87} & \cellcolor{gray!10}{43540} & \cellcolor{gray!10}{500} & \cellcolor{gray!10}{200} & \cellcolor{gray!10}{583}\\
Mujer & Adulto Maduro (36-43) & 99 & 43655 & 441 & 100 & 578\\
\cellcolor{gray!10}{Mujer} & \cellcolor{gray!10}{Mediana edad (44-53)} & \cellcolor{gray!10}{106} & \cellcolor{gray!10}{46240} & \cellcolor{gray!10}{436} & \cellcolor{gray!10}{150} & \cellcolor{gray!10}{496}\\
Mujer & Adulto Mayor (54-64) & 109 & 45280 & 415 & 120 & 526\\
\bottomrule
\end{tabular}}
\end{table}

\begin{figure}[H]

{\centering \includegraphics[width=1\linewidth,]{informe_files/figure-latex/obj5_correlacion-1} 

}

\caption{Correlación entre Volumen Operativo y Financiero}\label{fig:obj5_correlacion}
\end{figure}

En el gráfico de la figura \ref{fig:obj5_correlacion} se puede observar
la correlación que tiene el número transacciones con el total gastado
por su Segmento Demográfico. Se puede observar como la distribución
muestra una correlación positiva, principalmente impulsada por los
Segmentos Etarios de adultos jóvenes, esto significa que a mayor compra
hay mayor ganancia, esto puede resultar evidente pero es importante
resaltar que la mayoría de segmentos demográficos están por debajo de su
línea de correlación.

La curva de los hombres al estar más empinada que la de las mujeres,
indica que los hombres que compran con mayor frecuencia, su promedio de
compra es mayor que el de las mujeres que compran con la misma
frecuencia. Se observa también, que la línea morada empieza más alto que
la azul, lo que significa que las mujeres que compran con menor
frecuencia gastan más que los hombres que compran con menor frecuencia.
No obstante, a partir de la intersección que hay entre 100 compras y 105
compras, los hombres que compran con mayor frecuencia gastan más que las
mujeres que compran con mayor frecuencia. Analizando los extremos, se
puede apreciar como las mujeres jóvenes son las que mayor ganancia
generan de toda la distribución, siendo estas las que más compran y más
gastan, por otro lado, los hombres maduros son los que menos gasto
generan, seguido de los adultos jóvenes. Otro hallazgo interesante es
que las mujeres adultas generan más gasto con menos compras hechas que
las mujeres maduras. Y para concluir este análisis, se puede observar
como los adultos mayores de ambos segmentos demográficos son dos de los
segmentos que más compran, sin embargo, estas compras no son tan
elevadas como la de los adultos jóvenes.

\chapter{Conclusiones y Recomendaciones}

El presente estudio ha permitido caracterizar el comportamiento de
compra y la rentabilidad de los clientes en el sector retail mediante el
análisis descriptivo de un historial transaccional anual. A
continuación, se presentan las conclusiones derivadas de la evidencia
estadística y las recomendaciones estratégicas para la optimización de
los recursos comerciales.

\textbf{1.-} Se determinó que no existe una variación significativa en
el ticket promedio de compra entre hombres y mujeres. El análisis de
esperanza condicional nos permitió conocer que en ambos sectores se
presentaban valores casi idénticos (Mujeres: \$457; Hombres: \$455), con
medianas de \$120 y \$150 respectivamente y desviaciones estándar
comparables (566 y 554 respectivamente). Los diagramas de caja y bigotes
junto con los diagramas de densidad confirmaron distribuciones
fuertemente sesgadas a la derecha, con la mayoría de las transacciones
situándose por debajo de \$900.

Sin embargo, pudimos observar preferencias en las categorías de consumo,
siendo los hombres más predilectos a la categoría Electrónica, mientras
que las mujeres prefieren irse por Ropa, manteniéndose ambos en un nivel
similar de gasto en Belleza. Esto nos enseña que el género no nos sirve
para predecir la magnitud de gasto por transacción, pero si la
composición de la categoría de compra.

\textbf{Recomendación:} Se recomienda no dejar de lado ningún género en
ninguna categoria, ya que como casi todos gastan una cantidad
considerable en todas las cantegorias, descuidar la venta de alguna de
estas podría significar una pérdida.

\textbf{2.-} Se determinó que si existe una influencia de la edad y el
género en categorías especificas, por ejemplo, las mujeres mayores
tienen una mayor predilección por compras artículos de la categoría
``Ropa'', mientras que los hombres mayores se inclinan más a compras
artículos de la categoría ``Electrónica''. Esto sugiere que si hay una
predilección de la elección de un producto dependiendo de la edad y el
género.

\textbf{Recomendación:} Las campañas dirigidas a los hombres tienen que
ser orientadas a la Electrónica, especialmente a los hombres de mediana
edad y adultos mayores. Mientras que, las campañas dirigidas hacia las
mujeres, tienen que ser orientadas a la Ropa, especialmente a las
mujeres mayores de 36 años.

\textbf{3.-} Se refutó la hipótesis de que a mayor edad existe una mayor
preferencia por artículos ``Premium'' (\(> \$300\)). El segmento más
joven (18-26 años) lidera la proporción de compras de alto valor. Este
hallazgo nos permite saber la estrategia de negocio, ya que los adultos
jóvenes tienen mayor predilección por comprar artículos de alto costo.

\textbf{Recomendación:} Se recomienda hacer una estrategia del
posicionamiento de los productos de alto costo orientado a los adultos
jóvenes.

\textbf{4.-} Se determinó que el mayor gasto lo hacen los adultos
jóvenes, alcanzando un gasto acumulado superior a \textbackslash100K a
finales del ciclo del año 2023. No obstante, los otros Segmentos Etarios
también tienen un gasto considerable, por ejemplo, los de mediana edad y
los adultos mayores gastan \$90K en total del año 2023, indica que los
hay una estabilidad de consumo en la segunda mitad de vida.

\textbf{Recomendación:} Se puede hacer una un programa de financiamiento
para artículos de alta gama, para así atraer más público de los adultos
jóvenes. También, como el Segmento etario de 36-43 años son el que menos
gasta, podría hacerse una campaña de publicidad para que este grupo se
incentive a comprar más productos.

\textbf{5.-} Se concluye que existe una correlación positiva entre las
compras hechas y el total gastado. Según el análisis, los grupos que más
gasto y compras generan son los adultos jóvenes, pero también los
adultos de mayor edad generan una frecuencia de compra importante.

\textbf{Recomendación:} Como la mayoría de grupos compran mucho pero
gastan poco, se podría hacer algunas promociones o campañas de
publicidad para que estos grupos se incentiven a gastar más dinero del
que gastan normalmente.

%%%%%%%%%%%%%%%%%%%%%%%%%%%%%%%%%%%%%%%%%%%%%%%%%%%%%
%                   REFERENCIAS                     %
%%%%%%%%%%%%%%%%%%%%%%%%%%%%%%%%%%%%%%%%%%%%%%%%%%%%%
% Mantenemos la estructura natbib de tu plantilla
\renewcommand{\bibname}{Lista de Referencias}
\bibliographystyle{apalike} 
\bibliography{bibliografia.bib} 

%%%%%%%%%%%%%%%%%%%%%%%%%%%%%%%%%%%%%%%%%%%%%%%%%%%%%
%                   APÉNDICES                       %
%%%%%%%%%%%%%%%%%%%%%%%%%%%%%%%%%%%%%%%%%%%%%%%%%%%%%

\end{document}
